\chapter{Optimizer Architecture}

\section{Overview}

\begin{remark}
The previous chapters introduced the mathematical and physical
foundations of alignment modelling and optimization.
\end{remark}

\begin{remark}
In this chapter, these concepts are combined into a concrete system
architecture that enables the implementation of a full alignment
optimization engine.
\end{remark}

\begin{remark}
The architecture consists of the following main components:
\begin{itemize}
\item data ingestion (import and preprocessing),
\item initial guess construction,
\item constraint library,
\item objective engine,
\item numerical solver,
\item live visualization.
\end{itemize}
\end{remark}

---

\section{Data Flow}

\begin{remark}
The overall data flow can be summarized as:
\[
\text{Input Data}
\rightarrow
\text{Initial Guess}
\rightarrow
\text{Optimization}
\rightarrow
\text{Result}
\]
\end{remark}

\begin{remark}
In the implemented system, this flow is realized through a sequence of
well-defined transformations.
\end{remark}

\begin{verbatim}
Drop → Parser → Grabbeltisch → Initial Guess → SQP → Model → View
\end{verbatim}

---

\section{Optimization Session}

\begin{definition}[Optimization session]
An optimization session encapsulates all data and operations required
for a single optimization run.
\end{definition}

\begin{verbatim}
class OptimizationSession {
  constructor({ input, constraints, objectives }) {
    this.input = input
    this.constraints = constraints
    this.objectives = objectives
    this.initialGuess = null
    this.result = null
  }

  buildInitialGuess() {}
  runOptimization() {}
  commitToModel(projectModel) {}
}
\end{verbatim}

\begin{remark}
This abstraction separates:
\begin{itemize}
\item data preparation,
\item numerical optimization,
\item integration into the system model.
\end{itemize}
\end{remark}

---

\section{Initial Guess Construction}

\begin{remark}
The initial guess transforms raw geometric data into optimization
variables.
\end{remark}

\begin{verbatim}
function buildInitialGuess(polyline) {
  const s = computeArcLength(polyline)

  const theta = computeDirections(polyline)
  const kappa = diff(theta) / ds

  const kappaSmooth = smooth(kappa)

  const phi = kappaSmooth.map(k =>
    Math.asin(clamp(v*v*k/g, -1, 1))
  )

  return { kappa: kappaSmooth, phi }
}
\end{verbatim}

\begin{remark}
This step is essential for robust convergence of the optimization
algorithm.
\end{remark}

---

\section{Constraint Library}

\begin{definition}[Constraint]
A constraint is a function mapping the current state to a residual.
\end{definition}

\begin{verbatim}
function minRadius(kappa, Rmin) {
  return Math.abs(kappa) - 1/Rmin
}
\end{verbatim}

\begin{verbatim}
function parallelConstraint(p1, p2, d) {
  const dx = p1.x - p2.x
  const dy = p1.y - p2.y
  return Math.sqrt(dx*dx + dy*dy) - d
}
\end{verbatim}

\begin{remark}
All constraints are expressed uniformly as residuals and integrated into
the optimization problem.
\end{remark}

---

\section{Objective Engine}

\begin{definition}[Objective function]
The objective function is constructed as a weighted sum of residual
terms.
\end{definition}

\begin{verbatim}
function objectiveJerk(kappa, phi, ds, v, g) {
  const r = []
  for (let i = 0; i < kappa.length - 1; i++) {
    const dk = (kappa[i + 1] - kappa[i]) / ds
    const dphi = (phi[i + 1] - phi[i]) / ds
    const j = v*v*v*dk - g*v*Math.cos(phi[i]) * dphi
    r.push(j)
  }
  return r
}
\end{verbatim}

\begin{verbatim}
function buildObjectiveResiduals(ctx, objectives) {
  const r = []

  for (const obj of objectives) {
    const ri = obj.evaluate(ctx)
    const w = Math.sqrt(obj.weight ?? 1)
    for (const v of ri) r.push(w * v)
  }

  return r
}
\end{verbatim}

\begin{remark}
This formulation enables flexible adaptation to different design
criteria.
\end{remark}

---

\section{Residual Assembly}

\begin{remark}
The optimization problem is expressed entirely in terms of residuals.
\end{remark}

\begin{verbatim}
function buildResidualVector(ctx, constraints, objectives) {
  return [
    ...buildConstraintResiduals(ctx, constraints),
    ...buildObjectiveResiduals(ctx, objectives)
  ]
}
\end{verbatim}

---

\section{Numerical Solver}

\subsection{Gauss--Newton / SQP Core}

\begin{verbatim}
async function runSQPInteractive(x0, onStep) {
  let x = x0

  for (let k = 0; k < maxIter; k++) {

    const r = residuals(x)
    const J = jacobian(x)

    const p = solve(J, r)

    x = add(x, scale(p, alpha))

    onStep({ x, k, r })

    await nextFrame()
  }

  return x
}
\end{verbatim}

\begin{remark}
The solver operates on the residual formulation and exploits sparsity.
\end{remark}

---

\section{Live Visualization}

\begin{remark}
A key feature of the system is the live visualization of the optimization
process.
\end{remark}

\begin{verbatim}
optSession.runOptimization({
  onStep: ({ x, k }) => {
    const geom = reconstruct(x)

    geoRender.updateAlignment(geom)
    plotKappa(x.kappa)
    plotPhi(x.phi)
  }
})
\end{verbatim}

\begin{remark}
This enables interactive exploration and debugging of the optimization.
\end{remark}

---

\section{Network Extension}

\begin{remark}
The architecture extends naturally to networks of alignments.
\end{remark}

\begin{verbatim}
function residuals(network) {

  const r = []

  for (edge of network.edges) {
    r.push(...smoothness(edge))
  }

  for (node of network.nodes) {
    r.push(...nodeConstraints(node))
  }

  return r
}
\end{verbatim}

\begin{remark}
This allows simultaneous optimization of multiple connected tracks.
\end{remark}

---

\section{System Integration}

\begin{remark}
The optimization engine is integrated into the user workflow through the
interactive selection interface.
\end{remark}

\begin{verbatim}
onUserFinalizeSelection(selection) {

  const alignmentData =
    buildAlignmentFromSelection(selection)

  const optSession = new OptimizationSession({
    alignmentData
  })

  optSession.buildInitialGuess()
  optSession.runOptimization()
}
\end{verbatim}

---

\section{Summary}

\begin{summary}
The presented architecture combines mathematical modelling, numerical
optimization, and interactive visualization into a unified system.

By representing all constraints and objectives as residuals, the system
achieves a high degree of modularity and extensibility.

This architecture enables the transition from static alignment modelling
to dynamic, optimization-driven design of railway infrastructure.
\end{summary}
